\section{Completion of the canonical-index proof}\label{sec:index-zero}

\begin{proof}[Proof of Theorem~\ref{thm:fourfold-index-zero}]
Apply the canonical cover and the decomposition
\eqref{eq:index-product-cover} to \(X\).  Proposition~\ref{prop:index-product-case}
settles all decompositions with at least two positive dimensional
blocks and the purely abelian case.  The remaining case is a single
four-dimensional nonflat factor.  Corollary~\ref{cor:index-noncanonical}
settles its noncanonical quotients.  It remains to bound the orders
when \(X\) is canonical and has such a single-factor cover.

Suppose that these orders are unbounded, and choose a sequence whose
orders \(r\) tend to infinity.  Replace each member by its crepant
\(\Q\)-factorial terminalization \(V\), using
Lemma~\ref{lem:index-terminalization}.  The order \(r\) is unchanged.
The conclusions of Lemma~\ref{lem:index-single-factor-images} are
birational properties, so every proper positive dimensional rational
image of \(V\) has rationally connected resolution, and a smooth
projective model of \(V\) has irregularity zero.

By Lemma~\ref{lem:index-moving-divisor}, any member carrying an
effective nonbig divisor of positive Iitaka dimension has canonical
order dividing the fixed integer \(N_{\rm mov}\).  After discarding
a finite initial part of the sequence, every effective Weil divisor
of positive Iitaka dimension is therefore big.  For the cyclic
index cover \(\pi:Y\to V\), Lemma~\ref{lem:index-general-type-fibres}
now verifies the precise intermediate-fibration hypothesis of
Theorem~\ref{thm:scalar-bounds}, both for \(V\) with the trivial group
and for \(Y\) with its cyclic deck group.

Choose an ample rational Cartier divisor \(L\) on each \(V\), scaling
it by a positive rational number so that \(1\le L^4\le2\).
The scaling is possible without any uniform Cartier denominator.
The scalar estimates give constants \(C,c>0\), independent of the
member of the sequence, such that at very general points
\[
 \gamma(L;V)\le C2^{1/4},\qquad
 \varepsilon(L;x)\ge c,
 \qquad
 \varepsilon(\pi^*L;y)\ge c r^{1/4}.
\]
For the last inequality we used
\((\pi^*L)^4=rL^4\); the polarization upstairs is ample and invariant
under the deck group. Choose once and for all a rational number
$C_1\ge C2^{1/4}$ and put $b_0=1/(4(C_1+1))$.
Then $b_0(C_1+1)=1/4<3/4$ and
$3(2b_0)^4<3/16<1/4$, so both strict conditions in
\eqref{eq:frob-choice-b0} hold. Proposition~\ref{prop:quadratic-cost}
therefore gives, for every fixed integer $t\ge2$,
\[
 \varepsilon(P_t)\le 20(C_1+1)t^2.
\]
The constants are independent of the member of the sequence.
For each fixed product, its auxiliary choices are made before the
characteristic tends to infinity in that proposition. We then use
only finitely many such product sizes before letting the canonical
orders $r$ tend to infinity in the transport argument.

Lemma~\ref{lem:birational-countability} applies to the terminal
\(V\), since \(K_V\sim_{\Q}0\) and its smooth projective model has
irregularity zero.  Thus \(\Bir(V)\) is countable.  All hypotheses of
Proposition~\ref{prop:many-factor-contradiction} have now been checked:
the cover is a finite cyclic torsor on a big smooth open, the
polarizations have the two displayed lower bounds, the diagonal
quotients have the uniform quadratic upper bounds, and the
downstairs birational group is countable.  That proposition rules
out \(r\to\infty\), a contradiction.

The orders are consequently bounded in this final case as well.
There are only the cases just treated.  Taking a common multiple of
their bounded orders gives a single positive integer \(N_4\) with
\(N_4K_X\sim0\) for every input \(X\), as required.
\end{proof}
